
\subsection{Measures, Correction and Evaluation}
\label{sec:method_analysis}

\textbf{Measures.}
Angular error is the angle between the measured gaze direction
and the reference target. Following
\cite{niehorster2026fundamentals7}, we distinguish four properties:
\emph{accuracy}, the median angular error;
\emph{bias}, the median signed azimuth and elevation offsets;
\emph{dispersion} ($P_\mathrm{STD}$), the spread around the
mean gaze direction; and \emph{precision} (RMS-S2S),
the sample-to-sample jitter. For pursuit, dispersion is
computed from gaze--target residuals. We additionally report
P90, the 90th percentile of angular error. Evaluations are
performed per participant, with recording- and scene-level
aggregation to avoid domination by long sequences.
Stored wearer offsets were removed, so \emph{uncorrected}
denotes NeonNet estimates without person-specific correction.

\textbf{Offset correction.}
A constant offset is estimated from the median signed
azimuth and elevation errors of reference fixations.
Laboratory offsets are estimated from individual chessboards
(15 targets) or pooled recordings; in-situ offsets use
verified driving scenes. For each driving participant,
five folds reserve approximately 20\% of scenes for
calibration and evaluate corrections on the remaining
80\%. All methods use the same held-out scenes.
We additionally evaluate sparse calibration using
$k\in\{1,\ldots,8\}$ randomly or geometrically selected
references. Scenes, rather than individual gaze samples,
are split; calibration and test scenes may originate
from the same drive.

\textbf{Learned correction.}
To investigate whether systematic error dependencies can
be exploited, we compare the constant median offset with
linear and ridge regression models predicting signed
azimuth and elevation errors. Controlled models use
gaze position, pupil diameter and viewing distance,
individually or in combination. They are trained on
approximately 20\% of recordings and evaluated on
recording-disjoint data across five folds.

We then transfer corrections learned on chessboards to
held-out driving scenes using features available at
inference time. Finally, we compare median offsets with
regression models fitted directly on the driving calibration
scenes. Reference-target coordinates are used for training
and evaluation, not as prediction inputs.

\textbf{Statistics.}
Controlled comparisons use paired Wilcoxon tests,
Kruskal--Wallis tests across factor levels and Spearman
correlations for temporal trends, with Holm correction
for multiple comparisons. Driving corrections are
compared using paired scene-level tests with Holm
correction. We report uncertainty and fold-level
variability where appropriate.
